\section{Scope of the result}
\label{sec:scope}

The conclusion concerns one normalized Gram family with truncation
parameter $N$. Its arithmetic input is the divisibility relation and
the already known LCM operator. The extension established here is
the allowance of arbitrary admissible pointwise slowly varying
factors without multiplicativity, together with convergence throughout
the exact positive-exponent ideal range. The uniform spectral tails
and the consequences of Section~\ref{sec:consequences} are components
of that same result.

Several boundaries are worth stating explicitly. First, the hypotheses
do not merely require regular variation of
$\sum_{k\le N}|a(k)|^2$ for an arbitrary arithmetic sequence.
The pointwise representation $a(k)=k^{-\sigma}L(k)$ is what controls
the two rescaled factors in each off-diagonal Gram entry. The more
general multiplicative settings of~\cite{Hilberdink2017} therefore
remain distinct from the present theorem. Second, no convergence rate
uniform over all admissible slowly varying functions is asserted.
Third, the $N$-independent eigenvalue bound is proved for every
$\eta<\rho$, not at its endpoint. Fourth, no simultaneous
large-index relative singular-value asymptotic is claimed.

Finally, the LCM eigenvalue asymptotics, their spectral constant, and
the prime-based factorization used to obtain them belong to
\cite{HilberdinkPushnitski2023}. The comparison with its even-integer
Gram statement is bound to the accessed 2021 preprint version. The
source comparison here does not certify global priority or an open
problem in the later literature. No target Euler factors, functional
equation, prescribed zero set, or Hilbert--P\'olya realization is
identified by this Gram-limit theorem.
